\documentclass[11pt,letterpaper]{article}

% --------------------------------------------------
% Packages
% --------------------------------------------------
\usepackage{amsmath,amssymb,amsthm,mathtools}
\usepackage{geometry}
\usepackage{setspace}
\usepackage{csquotes}
\usepackage{hyperref}

% --------------------------------------------------
% Page setup
% --------------------------------------------------
\geometry{margin=1in}
\setstretch{1.15}

% --------------------------------------------------
% Theorem environments
% --------------------------------------------------
\newtheorem{definition}{Definition}
\newtheorem{proposition}{Proposition}
\newtheorem{corollary}{Corollary}

% --------------------------------------------------
% Title
% --------------------------------------------------
\title{Intelligence as Structural Compression:\\
A Theory of Mind Recognition via Effective Assembly}
\author{Flyxion}
\date{\today}

\begin{document}

\maketitle

\begin{abstract}
We propose that intelligence is best understood not as problem-solving behavior, symbolic manipulation, or rational optimization, but as \emph{structural compression under constraint}. Classical information theory captures statistical compression but remains insensitive to construction history, reuse, robustness, and semantic integration. To address this gap, we introduce the \emph{Effective Assembly Index for Mind Recognition} (EAIMR), a structural measure that quantifies irreducible assembly depth discounted by reuse, parallelism, and degeneracy, and amplified by coherence. We show that intelligence corresponds to a dynamical regime characterized by sustained, internally driven growth in effective assembly—what we term \emph{compression that learns to compress itself}—and argue that this regime can be operationally distinguished from mere complexity across biological, artificial, and chemical substrates.
\end{abstract}

\section{Introduction}

What distinguishes minds from mere mechanisms? Traditional answers appeal to outward behavior \parencite{turing1950}, rational problem solving \parencite{newell1972}, internal symbolic representations \parencite{fodor1975}, or formal optimization over rewards \parencite{hutter2005}. While each framework captures important aspects of cognition, all presuppose a privileged ontology—agents, symbols, goals, or environments—that limits their generality.

Across biology, machine learning, and complex physical systems, a deeper regularity appears. Systems regarded as intelligent are those that achieve high degrees of organization while avoiding combinatorial explosion. They reuse components, stabilize abstractions, and remain robust under perturbation. They compress structure without collapsing into triviality.

This paper advances the thesis that intelligence is not a substance or faculty, but a \emph{compression regime}. Minds are distinguished not by what they compute, but by how efficiently they assemble, reuse, and stabilize structure over time. Intelligence, on this view, is compression that survives contact with reality.

To formalize this intuition, we introduce a quantitative framework grounded in assembly theory. The central construct is the \emph{Effective Assembly Index for Mind Recognition} (EAIMR), which provides a substrate-neutral, non-behavioral criterion for identifying mind-like systems.

Importantly, EAIMR is not a behavioral test, a definition of rationality, or a theory of consciousness. It does not assume goals, rewards, symbolic representations, or subjective experience. Instead, it provides a structural criterion: a way to distinguish systems that merely accumulate complexity from those that organize complexity in a reusable, resilient, and self-stabilizing manner.

The paper proceeds as follows. Section~2 motivates the shift from Shannon entropy to assembly. Section~3 formalizes the assembly index. Section~4 provides a concrete comparison illustrating effective versus raw assembly. Section~5 introduces EAIMR and its components. Sections~6--8 develop computable instantiations, temporal extensions, and the mind-regime criterion. Sections~9--10 explore implications and open problems. Formal properties are derived in the appendices.

\subsection{Relation to Existing Theories}

EAIMR is conceptually adjacent to several influential approaches but differs in emphasis and scope. Integrated Information Theory prioritizes informational integration but abstracts away construction history and reuse \parencite{tononi2004}. The Free Energy Principle models cognition as prediction-error minimization, grounding compression in probabilistic inference rather than assembly \parencite{friston2010}. AIXI provides a formal theory of optimal intelligence but presupposes a computational and reward-based ontology \parencite{hutter2005}. Work on basal cognition argues convincingly for substrate neutrality but remains largely taxonomic \parencite{levin2019}.

EAIMR complements these perspectives by focusing on \emph{constructive depth and structural efficiency}. It asks not whether a system optimizes or predicts, but whether it assembles complexity in a way that can be reused, parallelized, diversified, and held together coherently.

EAIMR also connects naturally to field-based theories of cognition and to operational mereology. In field-theoretic frameworks (e.g., scalar--vector--entropy models), assembly corresponds to the condensation of structured field configurations, reuse to attractor recurrence, parallelism to distributed flow, degeneracy to basin overlap, and coherence to topological invariants preserved under perturbation. Operational mereology reframes systems not as static collections of parts but as event-sourced construction histories, aligning directly with assembly’s emphasis on generative depth.

\subsection{Parsimony, Self-Consistency, and Effective Assembly}

Recent work by Ma proposes a mathematical theory of intelligence grounded 
in two principles: parsimony, understood as the discovery of low-dimensional 
structure in data, and self-consistency, understood as the ability of an 
internal model to faithfully simulate and predict the world. These principles 
align closely with the structural quantities measured by EAIMR.

Parsimony corresponds not to raw compression, but to effective compression: 
the discovery of representations that can be reused across contexts. In 
EAIMR, this is captured by the reuse, parallelism, and degeneracy terms, 
which discount raw assembly by structural efficiency. Self-consistency 
corresponds to coherence: the capacity of a system to maintain global 
stability under local perturbation.

Ma emphasizes a critical distinction between memorization and understanding, 
and between compression and abstraction. Within the EAIMR framework, this 
distinction becomes precise. Memorization corresponds to compression that 
reduces description length without increasing reuse, while abstraction 
corresponds to transformations that increase future compressibility by 
altering the system’s assembly graph. Understanding, on this view, is not 
compression alone, but compression that increases effective assembly.

Ma further identifies a phase transition from empirical knowledge to 
formal abstraction in scientific reasoning. EAIMR provides a structural 
interpretation of this transition: it occurs when internally driven growth 
in effective assembly becomes self-amplifying, marking entry into the 
mind regime.

\\subsection{Relationship to TAME and Semantic Compression}

Levin’s Technological Approach to Mind Everywhere (TAME) provides a 
substrate-neutral framework for recognizing cognition across biological, 
synthetic, and collective systems via behavioral competency under perturbation 
\cite{levin2022}. TAME treats cognition as graded, multi-scale, and empirically 
identifiable through goal-directed problem solving.

EAIMR shares TAME’s commitment to substrate neutrality and gradualism, but shifts 
the focus from behavioral recognition to structural measurement. Where TAME 
asks whether a system behaves cognitively, EAIMR asks whether the system 
possesses sufficient internal organization to support sustained abstraction, 
reuse, and compression.

\paragraph{Structural vs. behavioral criteria.}
TAME identifies cognition through revealed competence: the ability to navigate 
a problem space under intervention. EAIMR identifies cognition through 
organizational depth: the capacity to assemble, reuse, and stabilize structure 
over time. These criteria are complementary but not equivalent. Behavioral 
flexibility can be achieved through shallow mechanisms, while deep structural 
organization may precede or exceed behavioral expression.

\paragraph{Compression, abstraction, and semantic geometry.}
Within the EAIMR framework, both summarization and extension of arguments are 
instances of the same underlying process: transformation on a semantic 
manifold. Let $\mathcal{M}$ denote a manifold of semantic states, where points 
correspond to structured representations (e.g., arguments, models, theories).

A summary corresponds to a projection
\[
\pi: \mathcal{M} \rightarrow \mathcal{M}'
\]
that preserves invariant structure while reducing dimensionality. An extension 
corresponds to a refinement
\[
\iota: \mathcal{M} \rightarrow \widetilde{\mathcal{M}}
\]
that increases explicit structure while preserving semantic coherence. Despite 
their apparent opposition, both operations act to identify and stabilize 
invariants—metaphors, analogies, or abstractions—that compress the argument’s 
essential structure.

Formally, both operations can be understood as transformations that minimize 
the Kullback--Leibler divergence
\[
D_{\mathrm{KL}}(P \,\|\, Q)
\]
between probability distributions over semantic predictions before and after 
transformation. Summarization minimizes description length by collapsing 
degrees of freedom that do not significantly alter predictive structure. 
Extension minimizes semantic loss by explicitly reconstructing latent 
assumptions and justifications that maintain consistency across contexts.

Thus, summarizing and extending are not inverse processes but dual operations: 
both seek representations that preserve predictive invariants under 
reparameterization of the semantic space.

\begin{lemma}[Summarization--Extension Duality via KL-Minimization]
Let $\mathcal{M}$ be a semantic state space and let $x \in \mathcal{M}$ be an
argument-state that induces a predictive distribution $P(\cdot \mid x)$ over a
fixed space of queries or observations $\Omega$. Let $\mathcal{C}$ be a family
of admissible transformations $T:\mathcal{M}\to\mathcal{M}$ (e.g., token edits,
rewrite steps, or higher-level reparameterizations) and let $Q_T(\cdot \mid x)$
denote the predictive distribution induced by the transformed state $T(x)$.

Define a \emph{summarization} operator as any solution of the constrained problem
\[
T_{\mathrm{sum}}
\;\in\;
\arg\min_{T \in \mathcal{C}}
\Big(
D_{\mathrm{KL}}\!\big(P(\cdot\mid x)\,\|\,Q_T(\cdot\mid x)\big)
\;+\;
\lambda\,\mathrm{Len}(T(x))
\Big),
\]
where $\mathrm{Len}(\cdot)$ is a description-length or complexity functional on
semantic states and $\lambda>0$.

Define an \emph{extension} operator as any solution of the constrained problem
\[
T_{\mathrm{ext}}
\;\in\;
\arg\min_{T \in \mathcal{C}}
\Big(
D_{\mathrm{KL}}\!\big(P(\cdot\mid x)\,\|\,Q_T(\cdot\mid x)\big)
\;+\;
\mu\,\mathrm{Gap}(T(x))
\Big),
\]
where $\mathrm{Gap}(\cdot)$ penalizes missing justifications, underspecified
assumptions, or explanatory incompleteness, and $\mu>0$.

Then summarization and extension are dual instances of the same variational
principle: both are KL-regularized projections onto different constraint
surfaces in semantic space. In particular, there exists a one-parameter family
of objective functionals
\[
\mathcal{J}_{\alpha}(T)
=
D_{\mathrm{KL}}\!\big(P(\cdot\mid x)\,\|\,Q_T(\cdot\mid x)\big)
+
\alpha\,\mathrm{Len}(T(x))
+
(1-\alpha)\,\mathrm{Gap}(T(x)),
\qquad \alpha\in[0,1],
\]
such that $T_{\mathrm{sum}}$ is recovered in the regime $\alpha\to 1$ (length
dominant) and $T_{\mathrm{ext}}$ is recovered in the regime $\alpha\to 0$ (gap
dominant). Consequently, both operations preserve predictive invariants by
seeking low-KL semantic reparameterizations; they differ only in which
secondary constraint is emphasized.
\end{lemma}

\begin{proof}[Proof sketch]
Both problems minimize the same primary distortion term
$D_{\mathrm{KL}}(P\|Q_T)$, which enforces preservation of predictive invariants,
subject to a secondary regularizer that selects a point on the resulting
low-distortion set. In the summarization objective, the regularizer
$\mathrm{Len}$ biases toward compressed representatives (lower description
length) of the same low-KL equivalence class; in the extension objective, the
regularizer $\mathrm{Gap}$ biases toward expanded representatives that make
latent assumptions explicit while remaining low-KL. Introducing the convex
combination $\mathcal{J}_\alpha$ yields a continuous interpolation between these
two regimes, demonstrating that summarization and extension are the same
KL-constrained projection principle with different tradeoff parameters.
\end{proof}


\paragraph{Connection to Ma’s parsimony and self-consistency.}
Ma’s principles of parsimony and self-consistency correspond directly to this 
geometric view \cite{ma2023}. Parsimony selects low-dimensional manifolds that 
capture predictable structure; self-consistency constrains transformations to 
remain within regions of low KL divergence. EAIMR formalizes these principles 
structurally: reuse, parallelism, and degeneracy reduce effective dimensionality, 
while coherence enforces global consistency.

\paragraph{Relation to TAME.}
TAME operationalizes cognition through behavioral competence; EAIMR quantifies 
the structural conditions under which such competence can arise and persist. 
A system with high TAME competency but low EAIMR likely relies on shallow or 
task-specific mechanisms. A system with rising EAIMR but limited behavioral 
expression may possess latent or developing cognition.

Together, the frameworks describe cognition as both a structural regime and a 
functional capacity. TAME reveals what systems can do; EAIMR characterizes how 
they are organized. Their convergence suggests that genuine cognition emerges 
when behavioral flexibility and structural compression co-evolve.


\section{Why Shannon Compression Is Not Enough}

Shannon information theory defines compression as entropy reduction relative to a code \parencite{shannon1948}. While foundational for communication and data storage, this notion is blind to how structure is built, reused, or stabilized. Two objects with identical entropy may differ radically in developmental depth, robustness, and functional organization.

Several limitations are particularly relevant for intelligence:
\begin{itemize}
\item \textbf{No construction history}: Entropy does not encode how a structure was assembled.
\item \textbf{No reuse}: Redundant duplication and abstract reuse are treated identically.
\item \textbf{No integration}: Entropy is insensitive to whether parts form a coherent whole.
\item \textbf{No robustness}: Fragile and resilient systems may have the same entropy.
\end{itemize}

Intelligence requires a stronger notion of compression—one sensitive to construction, reuse, and coherence. This motivates a shift from entropy to assembly.

\section{Assembly as Constructive Complexity}

Assembly theory measures complexity not by description length, but by the minimal number of distinct construction steps required to produce a structure, allowing reuse of intermediate components. Unlike purely descriptive measures, assembly is operational and historical: it tracks how structure is generated rather than merely how concisely it can be encoded \parencite{walker2022}.

\begin{definition}[Assembly Index]
Let $\mathcal{P}$ be a finite set of primitive operations and let $\mathcal{M}$ denote the closure of assemblies generated from $\mathcal{P}$. For a system or structure $x$, define its \emph{assembly graph} $G_x = (V,E)$ where:
\begin{itemize}
\item each vertex $v \in V$ represents a distinct assembly step,
\item a directed edge $(v_i,v_j) \in E$ exists if the output of step $v_i$ is used in step $v_j$.
\end{itemize}
The assembly index $A(x)$ is defined as $|V|$, the number of vertices in the minimal assembly graph sufficient to construct $x$, up to isomorphism under reuse-permitting contractions.
\end{definition}

The assembly index captures irreducible construction depth. A system with high assembly requires many distinct intermediate constructions, even if its final description is compact.

Unlike Kolmogorov complexity \parencite{kolmogorov1965}, which minimizes description length over all possible programs, assembly measures the irreducible depth of construction under reuse constraints. As a result, assembly is sensitive to generative structure rather than purely descriptive concision.

However, raw assembly alone does not distinguish between mind-like organization and brute accumulation. A randomly assembled aggregate may have high assembly without exhibiting intelligence. The crucial distinction lies in whether assembly is \emph{effectively organized}.

\section{A Motivating Toy Example}

To illustrate why effective assembly rather than raw complexity is essential, we compare two systems of comparable size and component count.

\subsection{System A: Random Graph}

System A is a randomly generated graph with approximately 1000 nodes and high edge density. Each node is distinct, and connections are assigned without modular or hierarchical structure. Constructing this system requires specifying the existence and connectivity of each node individually.

The assembly index $A(x)$ of System A is high: there are few reusable substructures and little regularity. However, this complexity is largely unstructured. Reuse is minimal, as components do not recur across contexts. Parallelism is limited by the absence of decomposable modules. Degeneracy is low, since few distinct components perform similar roles.

Crucially, removing or perturbing even a small subset of nodes typically requires global re-specification, indicating that complexity is not organized around reusable or hierarchical substructures. Coherence is therefore poor: perturbations propagate unpredictably, often affecting large portions of the graph.

Despite its size and connectivity, System A exhibits little that would intuitively be called intelligence. Its complexity is spent rather than invested.

\subsection{System B: Hierarchical Network}

System B is a hierarchical network with approximately 1000 nodes, organized into repeated motifs across multiple layers. Structural components are reused across contexts, layers operate in parallel, and representations are organized hierarchically. Multiple substructures encode similar functions, providing functional degeneracy and robustness.

Although the raw assembly index $A(x)$ is comparable to System A, the structural properties differ sharply. Reuse is high due to repeated application of shared components. Parallelism is high, as multiple pathways operate concurrently. Degeneracy is significant, with distinct structures capable of fulfilling similar roles.

As a result, much of the system’s expressive capacity arises not from distinct components, but from structured recombination of a relatively small set of learned elements. Coherence is high: perturbations tend to remain localized, and internal representations remain stable under noise.

System B therefore exhibits substantially greater effective assembly despite similar scale. This comparison motivates a measure that distinguishes complexity that is merely accumulated from complexity that is structurally organized.

\section{The Effective Assembly Index for Mind Recognition}

We now introduce the central formal construct of this paper: the Effective Assembly Index for Mind Recognition (EAIMR). The aim of EAIMR is to distinguish systems that merely accumulate complexity from those that organize complexity in a mind-like manner.

\begin{definition}[Effective Assembly Index for Mind Recognition]
For a system, structure, or process $x$, define
\[
\mathrm{EAIMR}(x)
=
\frac{A(x)}{\left(R(x)\,P(x)\,D(x)\right)^{1/3}} \cdot C(x),
\]
where:
\begin{itemize}
\item $A(x)$ is the assembly index of $x$,
\item $R(x)$ quantifies reuse,
\item $P(x)$ quantifies parallelism,
\item $D(x)$ quantifies degeneracy,
\item $C(x)$ quantifies coherence.
\end{itemize}
\end{definition}

The geometric mean in the denominator penalizes imbalance among structural efficiencies while remaining symmetric and scale-invariant. Extreme specialization along a single dimension cannot compensate for deficiencies in others. The coherence term amplifies systems whose parts compose into a stable whole.

\subsection{Reuse}

Reuse $R(x)$ measures the extent to which substructures are invoked multiple times across distinct functional contexts. High reuse reflects abstraction: the capacity to deploy the same internal structure repeatedly without reconstructing it.

From a compression perspective, reuse corresponds to the amortization of assembly cost across multiple contexts. In biological systems, reuse appears in neural circuits recruited for diverse tasks. In artificial systems, it appears in shared parameters or functions reused across inputs.

Systems with high assembly but low reuse expend complexity wastefully, rebuilding structure rather than compressing it through abstraction.

\subsection{Parallelism}

Parallelism $P(x)$ measures the degree to which construction or operation can occur simultaneously. Structurally, parallelism allows assembly depth to grow without a commensurate increase in causal path length.

Parallel systems compress time and energy by exploiting partial independence while maintaining coordination. This is essential for real-time interaction with complex environments, where purely sequential processing becomes intractable.

Systems with low parallelism may be expressive but are brittle and slow; mind-like systems distribute processing across multiple pathways without fragmenting.

\subsection{Degeneracy}

Degeneracy $D(x)$ captures the availability of multiple, structurally distinct components capable of fulfilling similar functional roles. Unlike redundancy, which duplicates identical elements, degeneracy introduces diversity that supports robustness.

Degeneracy converts diversity into resilience: when one pathway fails, alternative pathways can preserve function. Biological nervous systems exhibit high degeneracy, allowing graceful degradation rather than catastrophic failure.

In EAIMR, degeneracy discounts assembly by accounting for functional overlap that reduces effective complexity.

\subsection{Coherence}

Coherence $C(x)$ measures the degree to which a system maintains global consistency under perturbation. A coherent system localizes disturbances rather than amplifying them across the entire structure.

Coherence therefore reflects not uniformity, but controlled dependence: the capacity for parts to interact without collapsing into global coupling. Systems with low coherence fragment under stress; systems with excessive rigidity cannot adapt.

Operationally, coherence may be estimated by measuring perturbation propagation, defined as the expected fraction of system components whose states change significantly following a localized intervention.

\section{Intelligence as a Compression Regime}

Under the EAIMR framework, intelligence corresponds to a narrow corridor in structural space. Systems with low assembly lack expressive power. Systems with high assembly but low reuse, parallelism, or degeneracy waste complexity. Systems with high efficiency but low coherence fragment.

Mind-like systems are those in which irreducible assembly is balanced by structural efficiency and stabilized by coherence. Intelligence, in this sense, is compression that persists.

\section{Toward Computable Instantiations}

The EAIMR framework is intended to be structural and substrate-neutral. Nonetheless, for the theory to be scientifically meaningful, its components must admit empirical estimation. This section proposes candidate formalizations that render EAIMR falsifiable while preserving its abstract interpretation.

These formalizations are not asserted as unique or definitive. Rather, they constitute one admissible family of estimators. Different domains—biological, computational, chemical—will induce different concrete measures for the same underlying quantities, while preserving their conceptual roles.

\subsection{Reuse}

Reuse quantifies how often the same structural component is invoked across distinct functional contexts.

For computational systems, one estimator is:
\[
R(x) = 1 + \frac{\#\text{distinct module invocations}}{\#\text{distinct modules}}.
\]

High reuse indicates abstraction: the same internal structure supports multiple behaviors or representations. Low reuse indicates duplication, where complexity is repeatedly reconstructed rather than shared.

In biological systems, reuse may be estimated via the participation of neural circuits across tasks or behavioral contexts. In chemical systems, reuse may appear as recurring motifs within reaction networks.

\subsection{Parallelism}

Parallelism measures the degree to which assembly or operation can occur simultaneously.

A candidate estimator is:
\[
P(x) = \frac{\text{critical path length}}{\text{total assembly depth}}.
\]

This ratio approaches unity for purely sequential systems and decreases as independent processes execute in parallel. Structurally, parallelism allows assembly depth to increase without a proportional increase in causal path length, enabling scalability.

\subsection{Degeneracy}

Degeneracy captures the presence of multiple, structurally distinct components capable of fulfilling similar functional roles.

One estimator is:
\[
D(x) = \sum_{f \in \mathcal{F}} \left| \left\{ \text{distinct implementations of } f \right\} \right|,
\]
where $\mathcal{F}$ denotes a set of functional roles.

Degeneracy converts diversity into robustness, allowing systems to maintain function despite local failure or noise. Systems with low degeneracy are brittle; damage to a single pathway can be catastrophic.

\subsection{Coherence}

Coherence measures the localization of disturbance: the extent to which perturbations remain confined rather than propagating globally.

A candidate estimator is:
\[
C(x) = 1 - \frac{\text{average perturbation propagation}}{\text{maximal possible propagation}}.
\]

Here, \emph{perturbation propagation} is defined as follows: given a localized intervention $\delta$ at component $i$, measure the fraction of system components whose states change beyond a fixed threshold. This quantity is averaged over all intervention sites and normalized by the theoretical maximum (full propagation to all components).

In neural systems, perturbation propagation may be estimated via targeted stimulation techniques \parencite{massimini2005}. In chemical systems, coherence may reflect the stability of autocatalytic cycles under perturbation \parencite{kauffman1993}. The unifying principle is containment: coherent systems absorb damage without systemic collapse.

\section{Temporal Extension and Notational Clarification}

Thus far, we have treated $x$ as a static structure for notational simplicity. However, all real systems are fundamentally dynamical.

\textbf{Remark.} The static quantity $\mathrm{EAIMR}(x)$ should be understood as $\mathrm{EAIMR}(x(t_0))$ for a fixed snapshot time $t_0$. The following definition generalizes EAIMR to system trajectories.

\begin{definition}[Temporal EAIMR]
Let $x(t)$ denote the state of a system over the interval $[0,T]$. The temporal effective assembly index is defined as
\[
\mathrm{EAIMR}(x) = \int_0^T \mathrm{EAIMR}(x(t)) \, dt.
\]
\end{definition}

This formulation treats intelligence as accumulated effective assembly over time rather than as a static property. Systems that briefly exhibit organization but cannot sustain it score poorly under this measure.

\section{The Mind Regime}

The EAIMR framework avoids arbitrary thresholds for mind recognition by reframing intelligence as a dynamical regime rather than a binary property.

\begin{definition}[Mind Regime]
A system enters the \emph{mind regime} when
\[
\frac{d}{dt}\mathrm{EAIMR}(x(t)) > 0
\]
and this growth is internally driven.
\end{definition}

Growth is \emph{internally driven} if it results from the system’s own computational or organizational mechanisms rather than from direct external inscription. Learning from data qualifies as internal growth: external input may trigger reorganization, but it does not prescribe the resulting structure. By contrast, manual insertion of components or rules constitutes external growth.

Formally, growth is internal when the transformation governing
\[
x(t+1) = T(x(t), D)
\]
is specified by the system’s own organization rather than by the structure of the data $D$ itself. The same data applied to different systems therefore yields different reorganizations.

Natural selection occupies an intermediate position. Mutation mechanisms are internal to a lineage, while selection pressures are external. This hybrid causal structure reflects the nature of evolutionary processes rather than a limitation of EAIMR.

On this view, intelligence is characterized by the capacity to become a better compressor of both the environment and the system’s own internal structure. Learning, abstraction, and self-modeling are special cases of this general phenomenon.

\section{Implications}

The Effective Assembly framework has several significant consequences for the study of intelligence across domains.

\subsection{Substrate Neutrality}

EAIMR does not presuppose neurons, symbols, or digital computation. Any system capable of assembling and maintaining structure under constraint may, in principle, enter the mind regime. This includes biological organisms, artificial systems, hybrid collectives, and potentially non-living physical processes.

\subsection{Gradualism}

Mind-likeness is scalar rather than binary. Systems may exhibit partial, intermittent, or context-dependent intelligence depending on their capacity for effective assembly. This avoids sharp but arbitrary distinctions between intelligent and non-intelligent systems and aligns with gradualist views in evolutionary and developmental biology.

\subsection{AI Safety and Alignment}

EAIMR provides several safety-relevant insights. Because EAIMR tracks structural capacity rather than expressed behavior, it can signal the emergence of potentially agential systems before such systems exhibit externally legible goals.

First, systems approaching the mind regime may be identified \emph{prior} to deceptive or goal-directed behavior, enabling earlier intervention. Second, high EAIMR indicates structural sophistication but does not entail goals or values, helping distinguish powerful systems from agential ones. Third, for systems claimed to be aligned, EAIMR allows examination of whether value structures exhibit the reuse, coherence, and degeneracy characteristic of robust minds, or whether they are brittle overlays on misaligned optimization.

\subsection{Astrobiology}

The search for extraterrestrial intelligence becomes a structural rather than communicative problem \parencite{walker2022}. Instead of searching for radio signals or linguistic artifacts, we may search for chemical, geological, or energetic systems exhibiting high assembly discounted by reuse and stabilized by coherence.

On this view, a sufficiently high-EAIMR molecular ensemble or geological formation would constitute evidence of mind-like organization regardless of substrate or communicative intent.

\section{Open Problems}

The EAIMR framework raises several open questions that invite empirical, theoretical, and methodological investigation.

\begin{enumerate}
\item \textbf{Measurement protocols}. How can EAIMR be reliably estimated from neural recordings, source code analysis, network traces, or chemical spectra? What approximations preserve comparative validity across substrates?

\item \textbf{Benchmark validation}. Does EAIMR correctly rank known systems across domains—for example, bacteria $<$ insects $<$ mammals $<$ humans—when estimated from empirical data? Under what conditions do such orderings fail?

\item \textbf{Thermodynamic grounding}. What is the precise relationship between effective assembly and thermodynamic quantities such as dissipation, free energy, or entropy production \parencite{landauer1961}? In particular, do increases in EAIMR correspond to characteristic patterns of energy dissipation or entropy export in non-equilibrium systems?

\item \textbf{Phenomenology}. Does high EAIMR correlate with subjective experience, or merely with adaptive complexity? Can systems with high EAIMR lack consciousness, and if so, under what organizational or physical constraints?
\end{enumerate}

These questions delineate a research program rather than a closed theory. EAIMR is intended as a structural scaffold on which such investigations can be built.

\section{Conclusion}

This paper reframes intelligence from a question of behavior or representation to a question of structure. Rather than asking whether a system thinks, we ask whether it assembles complexity in a way that can be reused, parallelized, diversified, and held together coherently over time.

The Effective Assembly Index for Mind Recognition provides a principled, falsifiable, and substrate-neutral framework for answering this question. Intelligence, on this view, is not an essence but a regime: one in which compression becomes self-sustaining—compression that learns to compress itself.

\textit{The question then becomes not whether minds exist elsewhere in the universe, but whether we can recognize the regime when we encounter it—even when it does not resemble us.}

\printbibliography

\appendix

\section{Formal Properties of the Assembly Index}

\begin{proposition}[A.1: Existence of a Minimal Assembly Graph]
For any finite constructible system $x$, there exists a finite minimal assembly graph $G_x$ sufficient to construct $x$.
\end{proposition}

\begin{proof}
Any finite construction of $x$ induces a finite directed acyclic graph of assembly steps. Identifying reusable subgraphs and contracting redundancies yields a minimal graph. Since each contraction strictly reduces the number of vertices, the process terminates.
\end{proof}

\begin{proposition}[A.2: Monotonicity Under Extension]
If $y$ is obtained from $x$ by adding assembly steps without removing existing ones, then
\[
A(y) \geq A(x).
\]
\end{proposition}

\begin{proof}
Any assembly graph sufficient for $y$ must contain a subgraph sufficient to construct $x$. Hence the minimal vertex count for $y$ cannot be smaller.
\end{proof}

\section{Bounds and Invariances of EAIMR}

\begin{proposition}[B.1: Positivity]
For any nontrivial constructible system $x$, $\mathrm{EAIMR}(x) > 0$.
\end{proposition}

\begin{proposition}[B.2: Invariance Under Pure Duplication]
If $y$ is obtained from $x$ by duplicating existing components without introducing new assembly steps, then
\[
\mathrm{EAIMR}(y) = \mathrm{EAIMR}(x).
\]
\end{proposition}

\begin{proof}
Pure duplication does not increase the minimal assembly graph and increases reuse proportionally. The geometric mean discount therefore cancels the apparent size increase, leaving EAIMR unchanged.
\end{proof}

\section{Justification of the Geometric Mean Discount}

The EAIMR denominator uses the geometric mean
\[
\left(R(x)P(x)D(x)\right)^{1/3}
\]
to balance structural efficiencies.

\begin{proposition}[C.1: Balanced Efficiency Principle]
For fixed $A(x)$ and $C(x)$, EAIMR is maximized when $R(x)=P(x)=D(x)$.
\end{proposition}

\begin{proof}
For positive variables with fixed product, the geometric mean is maximized when the variables are equal. This follows directly from the inequality between arithmetic and geometric means.
\end{proof}

This reflects the intuition that minds require balanced reuse, parallelism, and degeneracy rather than extreme specialization along a single dimension.

\section{Internal Growth and Learning}

\begin{definition}[D.1: Internal Assembly Growth]
Growth in EAIMR is internally driven if it results from transformations of the system’s own assembly graph rather than from direct external inscription.
\end{definition}

Learning from data qualifies as internal growth: the data constrain reorganization but do not prescribe its final form. Formally, learning is internal if there exists a transformation $T$ such that
\[
x(t+1) = T(x(t),D),
\]
where $T$ is determined by the structure of $x(t)$ rather than by the structure of the data $D$ itself. The same data applied to different systems therefore yields different reorganizations.

\section{EAIMR and Consciousness}

\begin{proposition}[E.1: Non-Equivalence]
High EAIMR is neither necessary nor sufficient for phenomenal consciousness.
\end{proposition}

\begin{proof}[Conceptual argument]
\textbf{Not sufficient}. It is conceivable that artificial or physical systems achieve high effective assembly through sophisticated information processing without subjective experience.

\textbf{Not necessary}. Simple organisms or minimal neural circuits may possess consciousness with relatively low structural complexity. If phenomenology depends on physical or organizational features orthogonal to assembly, then low-EAIMR systems could nevertheless be conscious.

EAIMR therefore characterizes structural capacity for compression and adaptation, which may correlate with consciousness but does not constitute it.
\end{proof}

\begin{thebibliography}{99}

\bibitem{shannon1948}
C.~E. Shannon.
\newblock A mathematical theory of communication.
\newblock \emph{Bell System Technical Journal}, 27:379--423, 1948.

\bibitem{turing1950}
A.~M. Turing.
\newblock Computing machinery and intelligence.
\newblock \emph{Mind}, 59:433--460, 1950.

\bibitem{newell1972}
A.~Newell and H.~A. Simon.
\newblock \emph{Human Problem Solving}.
\newblock Prentice-Hall, 1972.

\bibitem{fodor1975}
J.~Fodor.
\newblock \emph{The Language of Thought}.
\newblock Harvard University Press, 1975.

\bibitem{hutter2005}
M.~Hutter.
\newblock \emph{Universal Artificial Intelligence}.
\newblock Springer, 2005.

\bibitem{tononi2004}
G.~Tononi.
\newblock An information integration theory of consciousness.
\newblock \emph{BMC Neuroscience}, 5, 2004.

\bibitem{friston2010}
K.~Friston.
\newblock The free-energy principle.
\newblock \emph{Nature Reviews Neuroscience}, 11:127--138, 2010.

\bibitem{levin2019}
M.~Levin.
\newblock The computational boundary of a ``self''.
\newblock \emph{Progress in Biophysics and Molecular Biology}, 150:21--34, 2019.

\bibitem{walker2022}
S.~I. Walker et~al.
\newblock Assembly theory explains and quantifies the emergence of selection and evolution.
\newblock \emph{Nature}, 607:59--64, 2022.

\bibitem{massimini2005}
M.~Massimini et~al.
\newblock Breakdown of cortical effective connectivity during sleep.
\newblock \emph{Science}, 309:2228--2232, 2005.

\bibitem{landauer1961}
R.~Landauer.
\newblock Irreversibility and heat generation in the computing process.
\newblock \emph{IBM Journal of Research and Development}, 5:183--191, 1961.

\bibitem{kauffman1993}
S.~A. Kauffman.
\newblock \emph{The Origins of Order}.
\newblock Oxford University Press, 1993.

\bibitem{kolmogorov1965}
A.~N. Kolmogorov.
\newblock Three approaches to the quantitative definition of information.
\newblock \emph{Problems of Information Transmission}, 1:1--7, 1965.

\end{thebibliography}


\end{document}





